% !TeX root = ../../Book2.tex
\section{Fourier 反演}
\label{b2:sec:0305}

反演的关键不是把指数的符号形式上改回来，而是说明两重振荡积分怎样得到原函数。直接交换两个没有绝对收敛保证的积分没有意义。我们先加一个可积的 Gauss 因子，使换序合法，再让这个因子逐渐趋于 \(1\)。

\subsection{光滑情形}
\label{b2:subsec:030501}

\begin{lemma}\label{b2:lem:03-gaussian-inversion}
设 \(d\) 为正整数，\(f\in L^1(\mathbb R^d;\mathbb C)\)、\(\varepsilon>0\)，并令 \(G_\varepsilon(y)=\varepsilon^{-d}e^{-\pi|y|^2/\varepsilon^2}\)（\(y\in\mathbb R^d\)）。则每个 \(x\in\mathbb R^d\) 满足
\begin{equation}\label{b2:eq:03-gaussian-inversion}
 \int_{\mathbb R^d}\widehat f(\xi)\cdot e^{-\pi\varepsilon^2|\xi|^2}e^{2\pi\mathrm{i}x\cdot\xi}\,\mathrm{d}\xi
 =(f*G_\varepsilon)(x).
\end{equation}
\end{lemma}
\begin{proof}
左侧绝对收敛，因为 \(\widehat f\) 有界而 Gauss 因子可积。代入变换定义，重积分绝对值的积分为 \(\|f\|_1\int e^{-\pi\varepsilon^2|\xi|^2}\,\mathrm{d}\xi<\infty\)，故可用 Fubini。内层正号积分由 Gauss 变换公式及偶性等于 \(G_\varepsilon(x-y)\)，从而得到右侧。
\end{proof}

为何先采用 Gauss 因子，可以与一维的硬截断作比较。对 \(R>0\) 定义
\[
 T_Rf(x)=\int_{-R}^R\widehat f(\xi)\cdot e^{2\pi\mathrm i x\xi}\,\mathrm d\xi.
\]
这个积分对 \(f\in L^1(\mathbb R)\) 有定义。代入变换后，绝对积分至多为 \(2R\|f\|_1\)，故交换次序合法，并得到
\[
 T_Rf(x)=\int_{\mathbb R}f(y)\cdot K_R(x-y)\,\mathrm dy,
 \quad K_R(z)=\frac{\sin(2\pi Rz)}{\pi z},\quad K_R(0)=2R.
\]
这里的卷积逐点绝对收敛，因为 \(|K_R(z)|\leq2R\)，故被积函数的绝对积分至多为 \(2R\|f\|_1\)。但是对每个 \(R>0\)，\(K_R\) 本身不属于 \(L^1(\mathbb R)\)。令 \(u=2\pi Rz\)，则
\[
 \int_{1/(2R)}^\infty|K_R(z)|\,\mathrm dz
 =\frac1\pi\int_\pi^\infty\frac{|\sin u|}{u}\,\mathrm du,
 \quad
 \int_{n\pi}^{(n+1)\pi}\frac{|\sin u|}{u}\,\mathrm du
 \geq\frac{2}{(n+1)\pi}\quad(n\geq1).
\]
最后一个下界使用每个半波的 \(|\sin u|\) 积分等于 \(2\)；把这些下界相加便得到发散的调和级数。因此硬截断核连单个核的绝对可积性都不满足，不能应用\cref{b2:thm:03-approximate-identity}直接推出 \(T_Rf\to f\)。Gauss 因子的逆积分则为非负核，并由上一节的归一化满足 \(\|G_\varepsilon\|_1=\int G_\varepsilon=1\)；其固定半径之外的质量随 \(\varepsilon\downarrow0\) 趋零，故满足逼近恒等的全部条件。两种方法在每个固定 \(R>0\) 或 \(\varepsilon>0\) 下都能合法换序，区别在于让 \(R\uparrow\infty\) 或 \(\varepsilon\downarrow0\) 时，是否已有可用的极限控制。

\ProseParagraphBreak

\Cref{b2:fig:03-hard-cutoff-gaussian-kernels}用一维的实际核比较这两种形状。取 \(R=2\)、\(\varepsilon=1/4\)，两核在原点的值都为 \(4\)，并使用同一组横、纵坐标刻度。硬截断核的正负旁瓣来自锐利的频率边界；Gauss 核始终非负，离开原点后迅速降低。有限画幅不能判断整条实线上的绝对积分是否有限，这一点仍由上面的半波下界和 Gauss 单位质量公式决定。

% 自拟函数图：同峰值4、同横纵尺度；sinc分两侧绘制并显式补入连续延拓值，避免零点商式。
\begin{figure}[htbp]
\centering
\begin{tikzpicture}[x=1.65cm,y=0.70cm,font=\small]
 \begin{scope}
  \draw[gray!60] (-1.5,-1) rectangle (1.5,4.2);
  \draw[gray!55] (-1.5,0)--(1.5,0);
  \foreach \xt in {-1.5,-1,-0.5,0,0.5,1,1.5}{
   \draw (\xt,-1)--(\xt,-0.90);
   \node[below] at (\xt,-1.35) {\(\xt\)};
  }
  \foreach \yt in {-1,0,1,2,3,4}{
   \draw (-1.5,\yt)--(-1.43,\yt);
   \node[left] at (-1.54,\yt) {\(\yt\)};
  }
  \draw[thick,AnalysisBlue] plot[domain=-1.5:-0.001,samples=181]
   (\x,{sin(720*\x)/(pi*\x)});
  \draw[thick,AnalysisBlue] plot[domain=0.001:1.5,samples=181]
   (\x,{sin(720*\x)/(pi*\x)});
  \draw[thick,AnalysisBlue] (-0.001,{sin(-0.72)/(-0.001*pi)})
   --(0,4)--(0.001,{sin(0.72)/(0.001*pi)});
  \node[anchor=south west] at (-1.5,4.4) {\(K_2(x)=\dfrac{\sin(4\pi x)}{\pi x}\)};
  \node[below right] at (1.5,-1.90) {\(x\)};
 \end{scope}
 \begin{scope}[xshift=6.2cm]
  \draw[gray!60] (-1.5,-1) rectangle (1.5,4.2);
  \draw[gray!55] (-1.5,0)--(1.5,0);
  \foreach \xt in {-1.5,-1,-0.5,0,0.5,1,1.5}{
   \draw (\xt,-1)--(\xt,-0.90);
   \node[below] at (\xt,-1.35) {\(\xt\)};
  }
  \foreach \yt in {-1,0,1,2,3,4}{
   \draw (-1.5,\yt)--(-1.43,\yt);
   \node[left] at (-1.54,\yt) {\(\yt\)};
  }
  \draw[thick,AnalysisBlue] plot[domain=-1.5:1.5,samples=201]
   (\x,{4*exp(-16*pi*\x*\x)});
  \node[anchor=south west] at (-1.5,4.4) {\(G_{1/4}(x)=4e^{-16\pi x^2}\)};
  \node[below right] at (1.5,-1.90) {\(x\)};
 \end{scope}
\end{tikzpicture}
\caption[硬频率截断核与 Gauss 核]{硬频率截断核与 Gauss 核。一维取 \(R=2\)、\(\varepsilon=1/4\)，两核在原点均为 \(4\)，均绘于 \(-1.5\leq x\leq1.5\)，纵坐标尺度相同。左图在原点采用连续延拓值，并有正负旁瓣；右图非负。图只显示有限范围内的形状，不代替正文对无穷远尾部的论证。}
\label{b2:fig:03-hard-cutoff-gaussian-kernels}
\end{figure}

\begin{proposition}\label{b2:prop:03-smooth-inversion}
设 \(d\) 为正整数。若 \(f\in C_c^\infty(\mathbb R^d;\mathbb C)\)，则 \(\widehat f\in L^1(\mathbb R^d;\mathbb C)\)，并有在 \(\mathbb R^d\) 上处处成立的绝对收敛反演公式 \(f=\check {\widehat f}\)。
\end{proposition}
\begin{proof}
取整数 \(m>d\)。对每个坐标 \(j\)，\(f\) 及中间导数 \(\partial_j^k f\)（\(1\leq k\leq m\)）都光滑、紧支集，因而可积。逐次应用\cref{b2:prop:03-derivative}，再在非零频率处选取绝对值最大的坐标，得到 \(|\widehat f(\xi)|\leq C|\xi|^{-m}\)。结合 \(|\widehat f|\leq\|f\|_1\)，便有全空间估计 \(|\widehat f(\xi)|\leq C'(1+|\xi|)^{-m}\)，于是变换可积。在\eqref{b2:eq:03-gaussian-inversion}左侧以 \(|\widehat f|\) 控制，得到 \(\check {\widehat f}(x)\)；右侧由有界一致连续函数的逼近恒等结论一致趋于 \(f\)。
\end{proof}

此证明只用到已有的 Gauss 积分、导数变换与逼近恒等，没有预设一般反演。因此它也提供了以后在测试函数空间上建立变换的起点。

\subsection{\texorpdfstring{\(L^1\)}{L¹} 反演}
\label{b2:subsec:030502}

\begin{theorem}[Fourier 反演]\label{b2:thm:03-inversion}
设 \(d\) 为正整数，\(f\in L^1(\mathbb R^d;\mathbb C)\) 且 \(\widehat f\in L^1(\mathbb R^d;\mathbb C)\)。对每个 \(x\in\mathbb R^d\)，绝对收敛积分
\[
 g(x)\coloneqq\int_{\mathbb R^d}\widehat f(\xi)\cdot e^{2\pi\mathrm{i}x\cdot\xi}\,\mathrm{d}\xi
\]
定义了 \(f\) 的唯一连续代表 \(g\in C_0(\mathbb R^d;\mathbb C)\)，即 \(g=f\) 几乎处处。若事先指定的有限值可测代表 \(f_0\colon\mathbb R^d\to\mathbb C\) 在某一点 \(x_0\in\mathbb R^d\) 连续，则 \(f_0(x_0)=g(x_0)\)。
\end{theorem}
\begin{proof}
首先，绝对收敛给出 \(|g(x)|\leq\|\widehat f\|_1\)，而 \(g(x)=\widehat{\widehat f}(-x)\) 使 Riemann--Lebesgue 引理可用于可积函数 \(\widehat f\)，从而 \(g\in C_0(\mathbb R^d)\)。由控制收敛，\eqref{b2:eq:03-gaussian-inversion}左侧一致趋于 \(g=\check {\widehat f}\)，因为误差不超过
\[
 \int|\widehat f(\xi)|\cdot|e^{-\pi\varepsilon^2|\xi|^2}-1|\,\mathrm{d}\xi\longrightarrow0.
\]
右侧由\cref{b2:thm:03-approximate-identity}在 \(L^1\) 中趋于 \(f\)。这与左侧的一致收敛如何合起来识别原函数，可以先在一个有限测度的球 \(B=B(0,N)\) 中比较：
\[
 \int_B|f-g|
 \leq\|f-G_\varepsilon*f\|_1
       +|B|\cdot\|G_\varepsilon*f-g\|_\infty
 \longrightarrow0.
\]
因此 \(f=g\) 在每个这样的球中几乎处处成立。让 \(N\) 遍历正整数，可数个零测例外集的并仍为零测集，便得到全空间的几乎处处等式。

若另有连续代表 \(h=f\) 几乎处处，则连续函数 \(h-g\) 几乎处处为零。如果它在某点不为零，连续性就给出一个球，使其绝对值在该球上有严格正的下界；球的 Lebesgue 测度为正，与几乎处处为零矛盾。这证明了连续代表的唯一性。

最后设 \(f_0\) 是已指定的代表，并在 \(x_0\) 连续。记 \(E=\{\,y\mid f_0(y)\ne g(y)\,\}\)，则 \(E\) 为零测集。每个球 \(B(x_0,1/n)\) 都有正测度，因此可选 \(y_n\in B(x_0,1/n)\setminus E\)。于是 \(y_n\to x_0\) 且 \(f_0(y_n)=g(y_n)\)。分别使用 \(f_0\) 在 \(x_0\) 的连续性和 \(g\) 的连续性，得到 \(f_0(x_0)=g(x_0)\)。
\end{proof}

这称为\term[fourier-fanyan-gongshi]{Fourier 反演公式}[Fourier inversion formula]。条件 \(\widehat f\in L^1\) 保证普通反演积分绝对收敛；它不是所有形式的反演都成立的必要条件。尤其对区间示性函数，此条件失败，但仍可讨论带正则化因子的反演。代表的区别也有实际内容：把零函数只在原点改成 \(1\)，变换仍恒为零，反演积分给出的仍是处处为零的连续代表，不会恢复在原点另行指定的值。

\subsection{点态反演}
\label{b2:subsec:030503}

沿用第一卷定义 14.4.1：设 \(d\) 为正整数，对 \(f\in L^1_{\mathrm{loc}}(\mathbb R^d;\mathbb C)\) 指定一个有限值可测代表 \(f_0\colon\mathbb R^d\to\mathbb C\)。若在 \(x\in\mathbb R^d\) 的小球内满足
\[
 \dfrac1{|B(x,r)|}\int_{B(x,r)}|f_0(y)-f_0(x)|\,\mathrm{d}y\longrightarrow0,
\]
则 \(x\) 为该代表的 Lebesgue 点。第一卷命题 14.4.6 证明了非负径向非增单位质量核在 Lebesgue 点处的逼近；本章以 \(f\in L^1\) 满足该命题对非紧支集核要求的全局可积性。Gauss 核和本章 Poisson 核均满足核的条件。

\begin{corollary}[Gauss 正则化反演]\label{b2:cor:03-pointwise-inversion}
设 \(d\) 为正整数，\(f\in L^1(\mathbb R^d;\mathbb C)\)，并指定一个有限值可测代表 \(f_0\colon\mathbb R^d\to\mathbb C\)。则在 \(f_0\) 的每个 Lebesgue 点 \(x\in\mathbb R^d\) 有
\[
 f_0(x)=\lim_{\varepsilon\downarrow0}
 \int\widehat f(\xi)\cdot e^{-\pi\varepsilon^2|\xi|^2}e^{2\pi\mathrm{i}x\cdot\xi}\,\mathrm{d}\xi.
\]
特别地，该结论几乎处处成立。
\end{corollary}
\begin{proof}
由\eqref{b2:eq:03-gaussian-inversion}，积分就是 \((G_\varepsilon*f_0)(x)\)。可以应用上述第一卷命题；为说明局部平均条件如何控制本章这个无限支集核，这里写出相应的估计。固定 \(f_0\) 的一个 Lebesgue 点 \(x\)，记
\[
 A(r)=\int_{|y|<r}|f_0(x-y)-f_0(x)|\,\mathrm dy,
 \quad g_\varepsilon(r)=\varepsilon^{-d}e^{-\pi r^2/\varepsilon^2}.
\]
给定 \(\eta>0\)，存在 \(\delta>0\)，使 \(A(r)\leq\eta\cdot\omega_dr^d\) 对 \(0<r\leq\delta\) 成立，其中 \(\omega_d\) 是单位球体积。在球内使用
\(g_\varepsilon(|y|)=g_\varepsilon(\delta)+\int_{|y|}^{\delta}(-g_\varepsilon'(r))\,\mathrm dr\)，由非负积分的 Tonelli 换序得
\[
 \begin{aligned}
 &\int_{|y|<\delta}G_\varepsilon(y)\cdot|f_0(x-y)-f_0(x)|\,\mathrm dy\\
 &\quad=g_\varepsilon(\delta)A(\delta)
       +\int_0^\delta(-g_\varepsilon'(r))A(r)\,\mathrm dr\\
 &\quad\leq\eta\left(g_\varepsilon(\delta)\cdot\omega_d\delta^d
       +\int_0^\delta(-g_\varepsilon'(r))\cdot\omega_dr^d\,\mathrm dr\right)\\
 &\quad=\eta\int_{|y|<\delta}G_\varepsilon(y)\,\mathrm dy\leq\eta.
 \end{aligned}
\]
括号中的最后等式把同一个分层积分用于常函数 \(1\)。球外含 \(f_0\) 和含常数 \(f_0(x)\) 的两部分，分别由
\[
 \varepsilon^{-d}e^{-\pi\delta^2/\varepsilon^2}\|f\|_1,
 \qquad |f_0(x)|\int_{|y|\geq\delta}G_\varepsilon(y)\,\mathrm dy
\]
控制；两者在 \(\varepsilon\downarrow0\) 时都趋零。先固定 \(\eta,\delta\)，取尺度极限，再令 \(\eta\downarrow0\)，就得 \((G_\varepsilon*f_0)(x)\to f_0(x)\)。最后用第一卷定理 14.4.2 的几乎处处 Lebesgue 点结论。
\end{proof}

这项结论不要求 \(\widehat f\) 可积，但不能把 Gauss 因子直接删去，也不能未经证明改成任意频率截断形状。一个有两个单侧极限的一维函数，在跳跃点经偶核正则化时趋向中点值：近处分别用左右极限，远处用核在离原点固定距离外的一致衰减估计。点本身的任意赋值不会被积分恢复。

\ProseParagraphBreak

例如对 \(f=\mathbf1_{(0,1)}\)，卷积可以完整算成
\[
 (G_\varepsilon*f)(x)
 =\int_{(x-1)/\varepsilon}^{x/\varepsilon}e^{-\pi z^2}\,\mathrm dz.
\]
当 \(0<x<1\) 时，两端分别趋于负、正无穷，极限为 \(1\)；当 \(x<0\) 或 \(x>1\) 时，积分区间进入同一侧尾部，极限为零。在 \(x=0,1\) 处，偶性使极限为 \(1/2\)。区间端点处的值不论补成什么，都不影响变换；正则化恢复的是左右平均。另一方面，原变换是指数相位乘 sinc 函数，其绝对积分发散，所以这个例子没有满足无正则化反演的假设。

\ProseParagraphBreak

反演还能确定上一节的一维 Poisson 变换。对 \(f(x)=e^{-2\pi t|x|}\)，直接计算已给出 \(\widehat f=P_t\)，且两者均可积。反演与偶性于是给出
\[
 \widehat {P_t}(\xi)=e^{-2\pi t|\xi|}\quad(d=1).
\]
两侧连续，所以最初的几乎处处等式提升为处处等式。

反演建立以后，还能证明与卷积定理互相对应的乘积公式。两个 \(L^1\) 函数的乘积未必可积；下面利用一个因子的可积 Fourier 变换，使这个因子有界，从而保证乘积可积。

\begin{proposition}\label{b2:prop:03-product-transform}
设 \(d\) 为正整数，\(f,g\in L^1(\mathbb R^d;\mathbb C)\)，且 \(\widehat g\in L^1(\mathbb R^d;\mathbb C)\)。选取反演给出的 \(g\) 的有界连续代表，则 \(fg\in L^1(\mathbb R^d;\mathbb C)\)，并有对每个 \(\xi\in\mathbb R^d\) 都成立的公式
\[
 \widehat{fg}(\xi)=\int_{\mathbb R^d}\widehat f(\xi-\eta)\cdot\widehat g(\eta)\,\mathrm d\eta.
\]
右侧的卷积积分绝对收敛。
\end{proposition}
\begin{proof}
反演给出 \(\|g\|_\infty\leq\|\widehat g\|_1\)，所以 \(\|fg\|_1\leq\|f\|_1\cdot\|\widehat g\|_1\)。将 \(g(x)=\int\widehat g(\eta)e^{2\pi\mathrm i x\cdot\eta}\,\mathrm d\eta\) 代入 \(fg\) 的变换。双重积分的绝对积分为 \(\|f\|_1\cdot\|\widehat g\|_1\)，故可用 Fubini 交换次序；内层对 \(x\) 的积分恰为 \(\widehat f(\xi-\eta)\)。反演最初若只几乎处处成立，也不改变乘积的等价类。最后用 \(\|\widehat f\|_\infty\leq\|f\|_1\) 保证所写卷积处处绝对收敛。
\end{proof}

乘积截断会在频率中产生卷积，空间卷积则变成频率乘积。这两个操作都是局部化的方法，但对频率的影响不同：把一个函数限在某一区域，通常会把其频率分散；用一个平滑核平均，则按该核的变换压低频率。后面讨论频率局部化时还会反复使用这两个公式的差别。

\subsection{唯一性与正则性的频率判据}
\label{b2:subsec:030504}

\begin{corollary}[唯一性]\label{b2:cor:03-uniqueness}
设 \(d\) 为正整数。若 \(f,g\in L^1(\mathbb R^d;\mathbb C)\) 且 \(\widehat f=\widehat g\)，则 \(f=g\) 在 \(\mathbb R^d\) 上几乎处处成立。
\end{corollary}
\begin{proof}
差 \(h=f-g\) 的变换恒为零，特别属于 \(L^1\)，故 Fourier 反演使 \(h=0\) 几乎处处。
\end{proof}

零函数在频率侧可积，所以对差函数的零变换使用反演，就能在全部 \(L^1\) 中得到唯一性。例如若 \(f*G_\varepsilon=0\) 对某个固定 \(\varepsilon>0\) 成立，则变换相乘后得到 \(\widehat f(\xi)e^{-\pi\varepsilon^2|\xi|^2}=0\)。Gauss 因子处处非零，因而 \(f=0\)。然而从带误差的卷积数据恢复 \(f\)，却要除以趋零的 Gauss 因子；唯一可恢复与稳定可恢复是不同问题。



\begin{proposition}\label{b2:prop:03-inverse-regularity}
设 \(d\) 为正整数，\(m\) 为非负整数，\(f\in L^1(\mathbb R^d;\mathbb C)\)，且 \(\int_{\mathbb R^d}(1+|\xi|)^m|\widehat f(\xi)|\,\mathrm{d}\xi<\infty\)。则 \(f\) 有一个代表 \(g\in C^m(\mathbb R^d;\mathbb C)\)，且对每个满足 \(|\alpha|\leq m\) 的多重指标 \(\alpha\in\mathbb N_0^d\)，都有 \(\partial^\alpha g\in C_0(\mathbb R^d;\mathbb C)\)。
\end{proposition}
\begin{proof}
零阶权已保证反演适用，取连续代表 \(g=\check {\widehat f}\)。每次对 \(x_j\) 求导，只给频率被积函数增加 \(2\pi\mathrm{i}\xi_j\)。所有相应差商都由假设中的可积矩控制，所以可逐阶求导，得到
\[
 \partial^\alpha g(x)=\int(2\pi\mathrm{i}\xi)^\alpha\widehat f(\xi)e^{2\pi\mathrm{i}x\cdot\xi}\,\mathrm{d}\xi.
\]
其绝对值由频率矩积分控制，连续性与无穷远消失由 Riemann--Lebesgue 引理给出。
\end{proof}

这与空间矩产生频率微分的命题构成对称关系。频率中的可积衰减可以恢复空间光滑性，条件必须作用于带权的绝对积分；仅知道变换逐点趋零，没有提供这些矩。反演因此不仅是一条恢复数值的公式，也是从频率估计提取正则性的工具。


\Cref{b2:tab:03-transform-rules}集中列出本章的运算规则，便于在后续方程计算中同时核对公式与假设。
% 本书自编查阅表；全部公式与条件来自第三章已证明结果。
\begin{table}[H]
\centering\small\renewcommand{\arraystretch}{1.2}\setlength{\tabcolsep}{4pt}
\caption[Fourier 变换的运算规则]{Fourier 变换的运算规则。\(d\) 为正整数，\(x,\xi\in\mathbb R^d\)，统一采用指数 \(e^{-2\pi\mathrm{i}x\cdot\xi}\)；允许复值函数，实值情形包含在内。第三列列出各行的可积性与参数条件。}
\label{b2:tab:03-transform-rules}
\begin{tabularx}{\linewidth}{@{}>{\raggedright\arraybackslash}p{3.4cm}>{\raggedright\arraybackslash}p{6.1cm}>{\raggedright\arraybackslash}X@{}}
\toprule
空间中的操作 & 频率中的表达 & 适用条件 \\
\midrule
\(f(x-h)\) & \(e^{-2\pi\mathrm{i}h\cdot\xi}\cdot\widehat f(\xi)\) & \(f\in L^1(\mathbb R^d)\)，\(h\in\mathbb R^d\) \\
\(e^{2\pi\mathrm{i}\eta\cdot x}\cdot f(x)\) & \(\widehat f(\xi-\eta)\) & \(f\in L^1(\mathbb R^d)\)，\(\eta\in\mathbb R^d\) \\
\(f(Ax)\) & \(|\det A|^{-1}\cdot\widehat f(A^{-\mathsf T}\xi)\) & \(f\in L^1(\mathbb R^d)\)，\(A\in\mathbb R^{d\times d}\)，\(\det A\ne0\) \\
\(\partial_jf\) & \(2\pi\mathrm{i}\xi_j\widehat f(\xi)\) & \(1\leq j\leq d\)，\(f\in L^1(\mathbb R^d)\)，弱导数 \(\partial_jf\in L^1(\mathbb R^d)\) \\
\(x_jf\) & \(\dfrac{\mathrm{i}}{2\pi}\partial_{\xi_j}\widehat f\) & \(1\leq j\leq d\)，\(f,x_jf\in L^1(\mathbb R^d)\) \\
\(f*g\) & \(\widehat f\cdot\widehat g\) & \(f,g\in L^1(\mathbb R^d)\) \\
\(\begin{gathered}G_\varepsilon(x)=\\\varepsilon^{-d}e^{-\pi|x|^2/\varepsilon^2}\end{gathered}\) & \(e^{-\pi\varepsilon^2|\xi|^2}\) & \(\varepsilon>0\)，此核可积且积分为 \(1\) \\
\bottomrule
\end{tabularx}
\end{table}

